\chapter{Tasks Performed}

\section{Dark Mode Implementation on Frontend Applications}

\clearpage

\section{Storybook Update for MCP}

\clearpage

\section{CI Report GCS Access}

\subsection{Project Overview}

\textbf{CI Report GCS Access} is a lightweight HTTP proxy developed to provide secure web access to Continuous Integration (CI) report artifacts stored in private Google Cloud Storage (GCS) buckets. The application exposes CI reports through standard HTTP endpoints while ensuring that the underlying storage remains inaccessible to the public.

The service is primarily intended to simplify the consultation of automatically generated reports, allowing developers and quality assurance teams to open reports directly from a web browser without requiring direct authentication to Google Cloud Storage.

The application is implemented using the \textbf{Hono} web framework and \textbf{TypeScript}, resulting in a lightweight, performant, and maintainable backend service.

\subsection{Objectives}

The project addresses several operational requirements encountered during the CI/CD workflow:

\begin{itemize}
    \item Keep Google Cloud Storage buckets private while allowing report visualization.
    \item Provide authenticated access through a simple HTTP interface.
    \item Eliminate the need for developers to manually interact with Google Cloud Storage.
    \item Simplify report sharing between development teams.
    \item Offer additional utilities such as report discovery and archive download.
\end{itemize}

\subsection{System Architecture}

The application acts as an intermediary between clients and Google Cloud Storage. Rather than exposing storage buckets publicly, every request passes through the proxy, which authenticates using a Google Cloud service account before retrieving the requested object.

The request processing workflow consists of the following steps:

\begin{enumerate}
    \item A client sends an HTTP request for a CI report.
    \item The proxy translates the request URL into the corresponding Google Cloud Storage object path.
    \item The requested object is retrieved using a service account with read-only permissions.
    \item The object is streamed directly back to the client.
\end{enumerate}

This architecture ensures that sensitive storage resources remain protected while maintaining a simple and transparent user experience.

\subsection{Technology Stack}

The application is built upon the technologies listed in Table~\ref{tab:gcs-tech-stack}.

\begin{table}[H]
\centering
\caption{Technology Stack}
\label{tab:gcs-tech-stack}
\begin{tabular}{|l|p{9cm}|}
\hline
\textbf{Technology} & \textbf{Purpose} \\
\hline
TypeScript & Main programming language \\
\hline
Hono & Lightweight HTTP framework \\
\hline
Node.js & JavaScript runtime environment \\
\hline
pnpm & Package manager \\
\hline
Google Cloud Storage & Storage of CI report artifacts \\
\hline
Google Cloud Service Account & Secure authentication to GCS \\
\hline
\end{tabular}
\end{table}

\subsection{Functionalities}

The service centralizes access to CI reports and provides several features that simplify report management and consultation.

\subsubsection{Secure Report Serving}

The primary functionality of the application is to retrieve report files stored in private GCS buckets and serve them securely over HTTP. HTML reports can therefore be opened directly from a browser without exposing the storage bucket publicly.

\subsubsection{Report Availability Verification}

Before attempting to open a report, users can verify whether a specific report exists. This operation avoids downloading unnecessary data and simply returns whether the requested resource is available.

\subsubsection{Automatic Report Discovery}

The application can search a bucket and identify the most recently uploaded report whose folder name matches a given prefix or suffix. This feature is particularly useful for locating the latest CI execution without knowing its exact identifier.

\subsubsection{Archive Download}

Entire report directories can be downloaded as compressed ZIP archives. Every file belonging to a report folder is included while preserving the directory structure.

\subsubsection{Bucket Listing}

The application can enumerate all Google Cloud Storage buckets accessible by the configured service account, providing administrators with a quick overview of available resources.

\subsubsection{Web User Interface}

Besides its REST API, the application offers a simple graphical interface allowing users to perform the available operations without manually constructing HTTP requests.

\subsection{Application Configuration}

The application requires a minimal configuration before execution.

The runtime environment is based on Node.js (version 24 or higher) together with the \texttt{pnpm} package manager (version 10 or higher). Authentication is performed using a Google Cloud service account possessing read-only permissions on the target storage bucket.

The main configuration parameters are summarized in Table~\ref{tab:gcs-env}.

\begin{table}[H]
\centering
\caption{Environment Variables}
\label{tab:gcs-env}
\begin{tabular}{|l|c|c|p{6cm}|}
\hline
\textbf{Variable} & \textbf{Required} & \textbf{Default} & \textbf{Description} \\
\hline
\makecell{GOOGLE\_CLOUD\\\_SERVICE\_\\CREDENTIALS\_JSON} & Yes & -- & JSON service account credentials used to authenticate with Google Cloud Storage. \\
\hline
NODE\_ENV & No & development & Runtime environment. \\
\hline
PORT & No & 3000 & HTTP server listening port. \\
\hline
\end{tabular}
\end{table}

\subsection{Application Programming Interface}

The application exposes several REST endpoints to interact with Google Cloud Storage.

\subsubsection{Report Proxy}

The \texttt{/gcs} endpoint serves report files stored inside a bucket. When a directory is requested, the application automatically returns the \texttt{index.html} file.

\begin{table}[H]
\centering
\caption{Report Proxy Endpoints}
\begin{tabular}{|p{6cm}|p{7cm}|}
\hline
\textbf{Endpoint} & \textbf{Description} \\
\hline
GET /gcs/\{bucket\}/\{reportId\}/ & Returns the report index page. \\
\hline
GET /gcs/\{bucket\}/\{reportId\}/\{objectPath\} & Returns a specific file from the report. \\
\hline
\end{tabular}
\end{table}

\subsubsection{Report Existence Verification}

The \texttt{/check-exist} endpoint verifies whether a report or one of its files exists inside a bucket.

\begin{table}[H]
\centering
\caption{Report Verification Endpoint}
\begin{tabular}{|p{6cm}|p{7cm}|}
\hline
\textbf{Endpoint} & \textbf{Purpose} \\
\hline
GET /check-exist/\{bucket\}/\{reportId\} & Returns whether the report exists. \\
\hline
\end{tabular}
\end{table}

\subsubsection{Report Search}

The \texttt{/search} endpoint retrieves the latest report folder matching either a prefix or a suffix.

\begin{table}[H]
\centering
\caption{Search Endpoints}
\begin{tabular}{|p{6cm}|p{7cm}|}
\hline
\textbf{Endpoint} & \textbf{Description} \\
\hline
GET /search/\{bucket\}/prefix=\{value\} & Finds the newest folder beginning with the specified prefix. \\
\hline
GET /search/\{bucket\}/suffix=\{value\} & Finds the newest folder ending with the specified suffix. \\
\hline
\end{tabular}
\end{table}

\subsubsection{Folder Download}

The \texttt{/download} endpoint generates a ZIP archive containing every file belonging to a report folder.

\begin{table}[H]
\centering
\caption{Download Endpoint}
\begin{tabular}{|p{6cm}|p{7cm}|}
\hline
\textbf{Endpoint} & \textbf{Description} \\
\hline
GET /download/\{bucket\}/\{folderPath\} & Downloads the selected report folder as a ZIP archive. \\
\hline
\end{tabular}
\end{table}

\subsubsection{Bucket Enumeration}

The application also exposes an endpoint for listing every Google Cloud Storage bucket accessible by the configured service account.

\begin{table}[H]
\centering
\caption{Bucket Listing Endpoint}
\begin{tabular}{|p{6cm}|p{7cm}|}
\hline
\textbf{Endpoint} & \textbf{Description} \\
\hline
GET /buckets & Returns all accessible GCS buckets. \\
\hline
\end{tabular}
\end{table}

\subsection{Testing}

The project includes integration tests that communicate directly with a dedicated Google Cloud Storage testing bucket. Authentication is performed using the same service account credentials employed by the application.

The tests validate the following functionalities:

\begin{itemize}
    \item Retrieval of report files.
    \item Report existence verification.
    \item Report search operations.
    \item ZIP archive generation.
    \item Bucket listing.
\end{itemize}

When the required credentials are not available, the integration tests are automatically skipped.

\subsection{Summary}

CI Report GCS Access provides a secure and efficient solution for exposing CI-generated reports stored in private Google Cloud Storage buckets. By acting as an authenticated proxy, the application preserves storage security while offering a simple HTTP interface for report visualization, search, verification, and download. Its lightweight architecture, modern technology stack, and integration with Google Cloud Storage make it an effective tool for improving accessibility to Continuous Integration artifacts within the development workflow.


\section{Visual Regression Testing Infrastructure}

\subsection{Introduction}

Modern web applications evolve continuously, and frequent modifications to the user interface can unintentionally introduce visual defects. Traditional software testing approaches, such as unit and integration testing, mainly verify application behavior, business logic, and component structure. However, they are not sufficient to detect presentation-level regressions affecting layout, typography, colors, images, or component positioning.

To address this limitation, a Visual Regression Testing (VRT) infrastructure was introduced using \textbf{Reg-Suit}. This system automatically captures application screenshots, compares them against approved reference images, detects visual differences, and generates reports that facilitate review during the Continuous Integration (CI) workflow.

The implemented solution covers two major application areas:

\begin{itemize}
    \item Storybook component visual testing.
    \item End-to-end (E2E) page visual testing.
\end{itemize}


\subsection{Visual Regression Testing}

\subsubsection{Concept}

Visual Regression Testing is a testing technique that detects unintended changes in the rendered appearance of a user interface by comparing screenshots generated from different application versions.

Unlike traditional automated tests that validate HTML structure, component properties, or application behavior, VRT operates at the presentation layer. Any modification affecting the visual output, such as spacing, typography, colors, icons, images, or element positioning, can be detected through pixel comparison.

\subsubsection{Testing Workflow}

The visual regression workflow consists of three main phases:

\begin{enumerate}
    \item \textbf{Snapshot generation:} screenshots are captured from the current application version.
    
    \item \textbf{Baseline comparison:} the generated screenshots are compared against reference snapshots stored from a previous validated version.
    
    \item \textbf{Difference analysis:} a pixel-level comparison algorithm identifies visual changes according to a predefined threshold and generates a regression report.
\end{enumerate}


\subsection{Reg-Suit Overview}

\textbf{Reg-Suit} is an open-source visual regression testing framework based on a plugin architecture. It manages the complete comparison lifecycle, including:

\begin{itemize}
    \item Collecting actual screenshots.
    \item Retrieving baseline screenshots.
    \item Performing pixel comparison.
    \item Generating HTML reports.
    \item Publishing visual regression artifacts.
\end{itemize}

The generated report contains an interactive visualization of detected differences and a JSON file describing the comparison results.

\subsection{System Configuration}

Reg-Suit is configured through a JSON configuration file. The main configuration parameters are shown below:

\begin{lstlisting}[caption={Reg-Suit configuration}]
{
  "core": {
    "workingDir": ".reg-e2e",
    "actualDir": "__e2e_tests_screenshots_/",
    "thresholdRate": 0.0001
  },
  "plugins": {
    "reg-simple-keygen-plugin": {
      "actualKey": "<ACTUAL_SHA>",
      "expectedKey": "<EXPECTED_SHA>"
    },
    "reg-publish-gcs-plugin": {
      "bucketName": "<BUCKET_NAME>",
      "pathPrefix": "<FOLDER_PATH>"
    }
  }
}
\end{lstlisting}


\subsubsection{Working Directory}

The \texttt{workingDir} parameter defines the temporary directory used by Reg-Suit during execution.

After completion, this directory contains:

\begin{itemize}
    \item HTML regression report.
    \item JSON comparison result.
    \item Actual screenshots.
    \item Expected screenshots.
    \item Generated image differences.
\end{itemize}

Separate working directories are maintained for Storybook and E2E tests to allow both reports to coexist.


\subsubsection{Actual Screenshot Directory}

The \texttt{actualDir} parameter specifies the location where freshly generated screenshots are stored.

This directory must correspond to the output path generated by the screenshot capture process:

\begin{itemize}
    \item \texttt{\_\_screenshots-playwright\_\_} for Storybook.
    \item \texttt{\_\_e2e\_tests\_screenshots\_\_} for E2E testing.
\end{itemize}


\subsubsection{Pixel Difference Threshold}

The \texttt{thresholdRate} defines the maximum acceptable pixel difference ratio.

The configured value:

\[
thresholdRate = 0.0001
\]

means that a regression is reported when more than 0.01\% of pixels differ between the actual and expected screenshots.

This strict threshold allows the detection of small visual changes while avoiding insignificant rendering variations.


\subsection{Snapshot Identification and Storage}

\subsubsection{Snapshot Keys}

Each visual regression execution is identified using two keys:

\begin{itemize}
    \item \textbf{Actual key:} identifies the screenshots generated by the current commit.
    \item \textbf{Expected key:} identifies the baseline screenshots used for comparison.
\end{itemize}

During Pull Request execution:

\begin{itemize}
    \item The actual key corresponds to the current commit SHA.
    \item The expected key corresponds to the target branch commit SHA.
\end{itemize}


\subsubsection{Storage Backend}

Visual regression artifacts are stored using a cloud storage backend.

The production CI environment uses:

\begin{itemize}
    \item Google Cloud Storage (GCS).
\end{itemize}

For local development, an S3-compatible storage solution is provided:

\begin{itemize}
    \item \textbf{s3rver} local object storage.
    \item Storage directory: \texttt{.s3-data/}.
    \item Server URL: \texttt{http://localhost:4568}.
\end{itemize}


\subsection{Visual Regression Commands}

The project provides scripts for both Storybook and E2E regression testing.

\begin{table}[H]
\centering
\caption{Visual Regression Commands}
\begin{tabular}{|p{5cm}|p{4cm}|p{4cm}|}
\hline
\textbf{Purpose} & \textbf{Storybook} & \textbf{E2E} \\
\hline
Capture screenshots &
pnpm book:screenshots:playwright &
pnpm test:visual:e2e
\\
\hline

Local execution &
pnpm reg:run:sb:local &
pnpm reg:run:e2e:local
\\
\hline

Remote execution &
pnpm reg:run:sb:remote &
pnpm reg:run:e2e:remote
\\
\hline

Open report &
pnpm reg:open:sb &
pnpm reg:open:e2e
\\
\hline
\end{tabular}
\end{table}


\subsection{Regression Report Generation}

After comparison, Reg-Suit generates an \texttt{out.json} file containing the complete comparison result.

Each screenshot belongs to one of four categories:

\begin{itemize}
    \item \textbf{Passed items:} no significant visual difference detected.
    \item \textbf{Failed items:} pixel difference exceeds the configured threshold.
    \item \textbf{New items:} screenshots existing only in the current version.
    \item \textbf{Deleted items:} screenshots removed from the current version.
\end{itemize}


\subsection{CI Pipeline Integration}

The visual regression system is integrated into the GitHub Actions workflow:

\texttt{.github/workflows/reg-vis-suit.yml}

The pipeline consists of four main jobs.

\subsubsection{Pipeline Jobs}

\paragraph{Setup Job}

The setup stage determines:

\begin{itemize}
    \item The current commit SHA.
    \item The target branch SHA.
    \item Whether corresponding snapshots already exist in GCS.
\end{itemize}


\paragraph{Expected Regression Job}

This job generates baseline screenshots from the expected commit.

It is skipped when the required baseline snapshots already exist.


\paragraph{Actual Regression Job}

This job captures screenshots from the current commit.

It runs in parallel with the expected regression job to reduce CI execution time.


\paragraph{Comparison Job}

The final stage:

\begin{itemize}
    \item Downloads expected and actual snapshots.
    \item Executes Reg-Suit comparison.
    \item Generates the visual regression report.
    \item Uploads the report to GCS.
    \item Publishes a Pull Request comment containing the detected changes.
\end{itemize}


\subsection{Handling Visual Regression Noise}

Visual regression systems are sensitive to any rendering variation. Therefore, several mechanisms were implemented to reduce false positive results.


\subsubsection{Disabling Animations}

Animations introduce unstable screenshots because their state may differ between executions.

To solve this problem, animations and transitions are disabled before capturing screenshots:

\begin{lstlisting}[]
*, *::before, *::after {
 animation: none !important;
 transition-duration: 0s !important;
}
\end{lstlisting}

This approach stabilizes components such as:

\begin{itemize}
    \item Loading animations.
    \item Progress indicators.
    \item Material UI Skeleton components.
    \item Animated SVG elements.
\end{itemize}


\subsubsection{Ignoring Non-Deterministic Components}

Some components cannot provide deterministic screenshots, such as:

\begin{itemize}
    \item Live counters.
    \item Random banners.
    \item External widgets.
\end{itemize}

A custom CSS utility class was introduced:

\begin{lstlisting}[language=HTML]
<div class="no-screenshot">
    ...
</div>
\end{lstlisting}

Elements using this class are hidden during screenshot generation to avoid unnecessary regressions.


\subsubsection{Waiting for Asset Loading}

Incomplete rendering of assets can generate false differences.

Before each screenshot, the system waits until:

\begin{itemize}
    \item Images are fully loaded.
    \item Fonts are available.
    \item Lazy-loaded resources are rendered.
\end{itemize}


\subsubsection{Mocking Dynamic Data}

Dynamic values such as dates, identifiers, and counters produce inconsistent screenshots.

To guarantee deterministic rendering:

\begin{itemize}
    \item E2E tests use predefined API fixtures.
    \item Mock Service Worker (MSW) intercepts network requests.
    \item Dynamic browser APIs are frozen in Storybook environments.
\end{itemize}

For Storybook, the following APIs are controlled:

\begin{itemize}
    \item \texttt{Date}.
    \item \texttt{Math.random}.
    \item \texttt{requestAnimationFrame}.
\end{itemize}


\subsection{Security Considerations}

Visual regression reports may contain internal application screenshots. Therefore, generated reports are stored in private cloud storage.

To provide controlled access without exposing the storage bucket publicly, a dedicated proxy service was developed:

\textbf{CI Report GCS Access}

This service authenticates with Google Cloud Storage and securely serves reports through HTTP.


\subsection{Conclusion}

The implemented Visual Regression Testing infrastructure provides an automated mechanism for detecting unintended UI changes during software development.

By combining Reg-Suit, Storybook, Playwright, Google Cloud Storage, and CI automation, the solution improves frontend reliability and reduces the risk of introducing visual regressions into production.

The integration of deterministic testing strategies, artifact management, and automated Pull Request reporting makes the system suitable for large-scale continuous integration environments.


\section{AI-Agents Driven Development}

\clearpage